% Actual class: 04
% Slide deck: L3 (full) Frequency slides
% Exact scope: pp. 1--32 (Parts 1--2)
% NTULearn supplement: two L3 video transcripts; test2_FT.m
% Human verified: Yes

\section{第 04 节课：Frequency-domain Image Enhancement}
\label{lecture:SC4061-L04}

\lecturemeta
  {SC4061-L04}
  {2026-09-01}
  {L3 (full) Frequency slides；\texttt{test2\_FT.m}}
  {L3 pp.~1--32（Parts 1--2）；pp.~33--45 留待 Part 3}
  {课程视频 L3 Parts 1--2（字幕核对）}
  {已确认（2026-09-01）}

\subsection{本讲主线}

Frequency-domain processing（频域处理）把 image（图像）分解为不同 spatial frequency（空间频率）和 orientation（方向）的 Fourier modes（傅里叶模态），修改各模态的 coefficients（系数），再合成为输出图像。它特别适合 global periodic disturbance（全局周期干扰）与按频率定义的 filtering；但 Fourier transform 始终可用不等于始终最方便，是否简化问题取决于目标 operator（算符）在 Fourier basis 中是否 diagonal 或 sparse。

\lecturetopic{SC4061-L04-T01}{Fourier Decomposition 与二维 Modes}

对 $M\times N$ image，采用 $x=0,\ldots,M-1$、$y=0,\ldots,N-1$，Fourier atom（傅里叶原子）为
\[
  \boxed{a_{u,v}(x,y)=
  \exp\!\left[j2\pi\left(\frac{ux}{M}+\frac{vy}{N}\right)\right]},
\]
其中 $u=0,\ldots,M-1$、$v=0,\ldots,N-1$ 是 mode indices（模态索引），不是脱离具体算符而独立存在的 eigenvalues（特征值）。$u/M$ 与 $v/N$ 表示 cycles per pixel（周期/像素）；相应 angular spatial frequencies（角空间频率）为 $2\pi u/M$ 与 $2\pi v/N$。

\begin{figure}[htbp]
  \centering
  \resizebox{0.94\textwidth}{!}{\begin{tikzpicture}[x=1cm,y=1cm,font=\small]
  \tikzset{tile/.style={draw,very thick,minimum width=2.5cm,minimum height=1.75cm}}

  \begin{scope}[shift={(0,2.3)}]
    \fill[gray!45] (0,0) rectangle (2.5,1.75);
    \draw[very thick] (0,0) rectangle (2.5,1.75);
    \node[above=1mm,align=center,font=\scriptsize] at (1.25,1.75)
      {$(u,v)=(0,0)$\\constant};
  \end{scope}

  \begin{scope}[shift={(3.4,2.3)}]
    \begin{scope}
      \clip (0,0) rectangle (2.5,1.75);
      \foreach \i in {0,...,9}{
        \pgfmathtruncatemacro{\shade}{20+65*mod(\i,2)}
        \fill[black!\shade] (0.25*\i,0) rectangle ++(0.25,1.75);
      }
    \end{scope}
    \draw[very thick] (0,0) rectangle (2.5,1.75);
    \node[above=1mm,align=center,font=\scriptsize] at (1.25,1.75)
      {larger $u$\\vertical stripes};
  \end{scope}

  \begin{scope}[shift={(0,0)}]
    \begin{scope}
      \clip (0,0) rectangle (2.5,1.75);
      \foreach \i in {0,...,6}{
        \pgfmathtruncatemacro{\shade}{20+65*mod(\i,2)}
        \fill[black!\shade] (0,0.25*\i) rectangle ++(2.5,0.25);
      }
    \end{scope}
    \draw[very thick] (0,0) rectangle (2.5,1.75);
    \node[below=1mm,align=center,font=\scriptsize] at (1.25,0)
      {larger $v$\\horizontal stripes};
  \end{scope}

  \begin{scope}[shift={(3.4,0)}]
    \fill[gray!15] (0,0) rectangle (2.5,1.75);
    \begin{scope}
      \clip (0,0) rectangle (2.5,1.75);
      \foreach \k in {-5,...,10}{
        \draw[line width=3.4pt,black!78] (-1+0.35*\k,0) -- ++(2.8,2.8);
      }
    \end{scope}
    \draw[very thick] (0,0) rectangle (2.5,1.75);
    \node[below=1mm,align=center,font=\scriptsize] at (1.25,0)
      {nonzero $u,v$\\oriented stripes};
  \end{scope}

  \draw[->,very thick,noteBlue] (6.5,2.0) -- (8.3,2.0)
    node[midway,above,align=center] {linear\ combination};
  \node[draw,very thick,rounded corners,fill=orange!12,
    minimum width=2.7cm,minimum height=2.4cm,align=center] at (10.0,2.0)
    {$f(x,y)$\\arbitrary image};
  \node[align=center,text=noteBlue] at (8.25,0.65)
    {$MN$ orthogonal modes\\for an $M\times N$ image};
\end{tikzpicture}
}
  \caption{二维 Fourier modes：$(u,v)$ 控制灰度变化的密度与方向，$MN$ 个正交 modes 构成 $M\times N$ image space 的一组完整 basis。}
  \label{fig:sc4061-l04-fourier-modes}
\end{figure}

Pixel basis（像素基）同样有 $MN$ 个 basis images：每个 basis image 仅一个 pixel 为 $1$。因此每个 pixel 是一个 coordinate（坐标），整个 image space 的维数为 $MN$；natural-image correlation（自然图像相关性）可能使真实数据集中于低维结构，但不改变 ambient vector-space dimension（环境向量空间维数）。

\lecturetopic{SC4061-L04-T02}{DFT、IDFT 与 Spectrum}

本讲采用 forward transform 不归一化、inverse transform 除以 $MN$ 的 convention：
\[
  \boxed{F(u,v)=\sum_{x=0}^{M-1}\sum_{y=0}^{N-1}
  f(x,y)\exp\!\left[-j2\pi\left(\frac{ux}{M}+\frac{vy}{N}\right)\right]},
\]
\[
  \boxed{f(x,y)=\frac1{MN}\sum_{u=0}^{M-1}\sum_{v=0}^{N-1}
  F(u,v)\exp\!\left[j2\pi\left(\frac{ux}{M}+\frac{vy}{N}\right)\right]}.
\]
Normalization（归一化）也可对称分配为正、逆变换各 $1/\sqrt{MN}$；不同 convention 只缩放 coefficients，不改变可逆性与频率内容。

正交性的关键是一维 geometric sum（等比和）：
\[
  \sum_{x=0}^{M-1}e^{j2\pi x(u-s)/M}
  =\begin{cases}M,&u=s\pmod M,\\0,&u\ne s\pmod M.\end{cases}
\]
两个方向相乘后，仅相同的 $(u,v)$ mode 被保留，因此 DFT 给出唯一 coefficients。

\begin{center}
\small
\begin{tabularx}{\textwidth}{@{}>{\raggedright\arraybackslash}p{3.2cm}X@{}}
  \toprule
  Quantity（量） & Interpretation（解释） \\
  \midrule
  $F(0,0)$ & DC component（直流分量），$F(0,0)=\sum_{x,y}f(x,y)=MN\,\bar f$ \\
  $|F(u,v)|$ & magnitude spectrum（幅度频谱），表示该 mode 的强度 \\
  $\angle F(u,v)$ & phase spectrum（相位频谱），保存 spatial shift 与结构位置 \\
  $F(M-u,N-v)=F^*(u,v)$ & real-valued image 的 conjugate symmetry（共轭对称性） \\
  \bottomrule
\end{tabularx}
\end{center}

原始 DFT 把 DC 放在数组角点；\texttt{fftshift} 只重新排列 indices，使低频位于中心、离中心越远频率越高。IDFT 前必须 \texttt{ifftshift}。显示 spectrum 时宜使用 $\log(1+|F|)$，避免 DC 的巨大 dynamic range 掩盖其他 coefficients。

\textbf{对应例题：}
\relatedexercise{SC4061-L04-E01}{Constant Image 的 DFT 与 DC Component}。

\lecturetopic{SC4061-L04-T03}{FFT 与 Convolution Theorem}

FFT（Fast Fourier Transform，快速傅里叶变换）是高效计算 DFT 的算法，不是新的 transform。Direct 2D DFT 的朴素复杂度为 $O(M^2N^2)$；separable 2D FFT 约为
\[
  O\!\left(MN(\log M+\log N)\right).
\]

Convolution theorem（卷积定理）给出
\[
  g=f*h
  \quad\Longleftrightarrow\quad
  \boxed{G(u,v)=F(u,v)H(u,v)}.
\]
DFT 默认对应 circular convolution（循环卷积）；要获得 ordinary linear convolution，大小分别为 $M_1\times N_1$ 与 $M_2\times N_2$ 的输入至少需 zero-pad 到
\[
  (M_1+M_2-1)\times(N_1+N_2-1).
\]

若 kernel 为 $P\times Q$，两种方法的主要复杂度约为
\[
  \text{direct spatial convolution}: O(MNPQ),\qquad
  \text{FFT method}: O(MN\log(MN)).
\]
FFT 还需 transforms、padding 与 complex arithmetic。Small local kernel（如 $3\times3$）通常直接在 spatial domain 更快；large shift-invariant kernel 或直接按 frequency 设计的 global filter 更适合 FFT。Gaussian kernel 还可用 separability（可分离性）在 spatial domain 高效实现。

\lecturetopic{SC4061-L04-T04}{Frequency-domain Filtering Framework}

\begin{figure}[htbp]
  \centering
  \resizebox{0.99\textwidth}{!}{\begin{tikzpicture}[node distance=0.65cm and 0.7cm,>=Latex,font=\small]
  \tikzset{
    data/.style={draw,rounded corners,very thick,fill=blue!8,
      minimum width=2.0cm,minimum height=1.0cm,align=center},
    op/.style={draw,circle,very thick,fill=orange!15,minimum size=1.05cm,align=center}
  }

  \node[data] (f) {$f(x,y)$\\spatial image};
  \node[op,right=of f] (fft) {DFT};
  \node[data,right=of fft] (F) {$F(u,v)$\\spectrum};
  \node[op,right=of F] (mul) {$\times$};
  \node[data,above=0.75cm of mul] (H) {$H(u,v)$\\filter mask};
  \node[data,right=of mul] (G) {$G=HF$\\filtered spectrum};
  \node[op,right=of G] (ifft) {IDFT};
  \node[data,right=of ifft] (g) {$g(x,y)$\\output image};

  \draw[->,very thick,noteBlue] (f) -- (fft);
  \draw[->,very thick,noteBlue] (fft) -- (F);
  \draw[->,very thick,noteBlue] (F) -- (mul);
  \draw[->,very thick,noteBlue] (H) -- (mul);
  \draw[->,very thick,noteBlue] (mul) -- (G);
  \draw[->,very thick,noteBlue] (G) -- (ifft);
  \draw[->,very thick,noteBlue] (ifft) -- (g);

  \node[below=0.58cm of fft,align=center,text=gray!75!black]
    {optional \texttt{fftshift}\\centre DC for display};
  \node[below=0.58cm of ifft,align=center,text=gray!75!black]
    {undo with \texttt{ifftshift}\\before IDFT};
\end{tikzpicture}
}
  \caption{Frequency-domain filtering：把 spectrum 与 transfer function $H(u,v)$ 逐点相乘，再经 IDFT 返回空间域。}
  \label{fig:sc4061-l04-frequency-pipeline}
\end{figure}

完整流程为
\[
  F=\mathcal F\{f\},\qquad
  G=HF,\qquad
  g=\mathcal F^{-1}\{G\}.
\]
$H=1$ 表示完全保留，$H=0$ 表示移除，$0<H<1$ 表示衰减。Low frequency（低频）通常对应 smooth trend、illumination 与大尺度结构；high frequency（高频）通常对应 edge、fine texture 与 abrupt change。不能把 ``低频''直接等同于 signal、``高频''直接等同于 noise：white noise 具有宽频谱，periodic noise 可能只形成少数尖峰，而有效 edges 本身也是高频 signal。

\lecturetopic{SC4061-L04-T05}{Low-pass、High-pass 与 Notch Filtering}

令 centred spectrum 中
\[
  D(u,v)=\sqrt{u^2+v^2}
\]
表示到 DC 的距离。Ideal low-pass filter（理想低通）为
\[
  H_{\mathrm{ILPF}}(u,v)=
  \begin{cases}1,&D(u,v)\le D_0,\\0,&D(u,v)>D_0.\end{cases}
\]
其 sharp cutoff（突变截止）在空间域对应带 sidelobes 的 sinc-like kernel，容易在强边缘产生 ringing artifacts（振铃伪影）。讲义中若写成 $u^2+v^2\le D_0$，当 $D_0$ 表示 radius 时，严格形式应为 $u^2+v^2\le D_0^2$。

Gaussian low-pass filter（高斯低通）平滑衰减：
\[
  \boxed{H_{\mathrm{GLPF}}(u,v)=
  \exp\!\left[-\frac{D^2(u,v)}{2\sigma_f^2}\right]}.
\]
$\sigma_f$ 越大，frequency passband 越宽，图像越清晰；对应 spatial Gaussian kernel 越窄。两个 domains 的宽度成反比，Gaussian 没有 ideal filter 的明显 sidelobes，因此 ringing 更弱。

High-pass filter（高通）可由
\[
  H_{\mathrm{HP}}=1-H_{\mathrm{LP}}
\]
获得，用于提取 edge 与 detail。直接 high-pass 输出主要是细节；sharpening（锐化）通常需要把细节加回原图：$g_{\mathrm{sharp}}=f+\alpha g_{\mathrm{HP}}$。

\begin{figure}[htbp]
  \centering
  \resizebox{0.96\textwidth}{!}{\begin{tikzpicture}[x=1cm,y=1cm,font=\small]
  \tikzset{mask/.style={draw,very thick,minimum width=2.3cm,minimum height=2.3cm}}

  \begin{scope}[shift={(0,0)}]
    \fill[black!82] (0,0) rectangle (2.3,2.3);
    \fill[white] (1.15,1.15) circle (0.62);
    \draw[very thick] (0,0) rectangle (2.3,2.3);
    \draw[dashed,gray] (1.15,0) -- (1.15,2.3);
    \draw[dashed,gray] (0,1.15) -- (2.3,1.15);
    \node[below] at (1.15,-0.18) {Ideal low-pass};
  \end{scope}

  \begin{scope}[shift={(3.2,0)}]
    \fill[black!85] (0,0) rectangle (2.3,2.3);
    \fill[black!55] (1.15,1.15) circle (0.93);
    \fill[black!30] (1.15,1.15) circle (0.69);
    \fill[black!10] (1.15,1.15) circle (0.40);
    \draw[very thick] (0,0) rectangle (2.3,2.3);
    \draw[dashed,gray] (1.15,0) -- (1.15,2.3);
    \draw[dashed,gray] (0,1.15) -- (2.3,1.15);
    \node[below] at (1.15,-0.18) {Gaussian low-pass};
  \end{scope}

  \begin{scope}[shift={(6.4,0)}]
    \fill[white] (0,0) rectangle (2.3,2.3);
    \fill[black!82] (1.15,1.15) circle (0.62);
    \draw[very thick] (0,0) rectangle (2.3,2.3);
    \draw[dashed,gray] (1.15,0) -- (1.15,2.3);
    \draw[dashed,gray] (0,1.15) -- (2.3,1.15);
    \node[below] at (1.15,-0.18) {Ideal high-pass};
  \end{scope}

  \begin{scope}[shift={(9.6,0)}]
    \fill[white] (0,0) rectangle (2.3,2.3);
    \fill[black!82] (0.70,1.55) circle (0.20);
    \fill[black!82] (1.60,0.75) circle (0.20);
    \draw[very thick] (0,0) rectangle (2.3,2.3);
    \draw[dashed,gray] (1.15,0) -- (1.15,2.3);
    \draw[dashed,gray] (0,1.15) -- (2.3,1.15);
    \node[below] at (1.15,-0.18) {Notch reject pair};
  \end{scope}

  \draw[->,very thick,noteBlue] (5.35,3.0) -- (6.55,3.0)
    node[midway,above] {distance $D(u,v)$};
  \node[anchor=east,text=noteBlue] at (5.25,3.0) {DC / low frequency at centre};
\end{tikzpicture}
}
  \caption{以 centred spectrum 表示的常见 masks。白色表示 pass，黑色表示 reject；notch pair 同时抑制实值图像中共轭对称的周期噪声峰。}
  \label{fig:sc4061-l04-filter-masks}
\end{figure}

Notch reject filter（陷波滤波器）只抑制特定 frequencies，适合 spectrum 中呈 conjugate peak pair（共轭峰对）的 sinusoidal interference。Band-reject filter（带阻滤波器）删除一段 frequency band；若干扰只占少数尖峰，notch 往往比删除整个 band 更能保留图像信息。

\textbf{对应例题：}
\begin{itemize}
  \item \relatedexercise{SC4061-L04-E02}{Gaussian Low-pass 的 $\sigma_f$}；
  \item \relatedexercise{SC4061-L04-E03}{Periodic Noise 与 Notch Pair}。
\end{itemize}

\lecturetopic{SC4061-L04-T06}{何时 Fourier Basis 会简化问题}

Fourier modes 可在不指定算符时作为一组 periodic orthogonal basis（周期正交基）定义；$u,v$ 首先只是 mode labels。它们是 cyclic translations（循环平移）的共同 eigenvectors，也使所有 periodic linear shift-invariant operators（周期线性平移不变算符），包括 circular convolution，变为 diagonal。不同 operators 共享 eigenvectors，但 eigenvalues 不同：shift 的 eigenvalue 为 $e^{-j2\pi u/M}$，convolution 的 eigenvalue 为 kernel spectrum $H(u,v)$。

并非所有 operators 都被 Fourier transform 简化。例如 spatially varying pointwise multiplication（空间变化的逐点乘法）
\[
  g(x,y)=m(x,y)f(x,y)
\]
在 pixel basis 中已经 diagonal，而在 frequency domain 中变成 convolution $G=M*F$，会混合多个 modes。应选择使目标 operator diagonal 或 sparse 的 domain：小型 local/spatially varying 操作常在 spatial domain 更自然；large shift-invariant convolution、global periodic interference 与 frequency-selective filtering 常在 frequency domain 更自然。

\subsection{MATLAB 速查}

附件代码的核心对应关系为
\begin{verbatim}
F = fftshift(fft2(f));
G = H .* F;
g = real(ifft2(ifftshift(G)));
imagesc(log1p(abs(F)));
\end{verbatim}
直接保留 centred spectrum 的一个矩形区域相当于 crude ideal low-pass（粗糙理想低通），可能产生方向偏置与 ringing；更平滑的 Gaussian mask 通常更稳健。

\subsection{一分钟回顾}

DFT 用 $MN$ 个二维 complex exponential modes 表示 $M\times N$ image；magnitude 表示 mode 强度，phase 保存空间位置，\texttt{fftshift} 只改变排列。FFT 高效计算 DFT，convolution theorem 使 shift-invariant convolution 变成 spectrum multiplication，但小 kernel 与 spatially varying operator 不一定适合 FFT。Low-pass 保留趋势、high-pass 保留细节，Gaussian cutoff 比 ideal cutoff 更少 ringing；periodic noise 常以共轭尖峰出现，应成对使用 notch reject。
